\section{Ablation study}\label{sec:ablation}

This section examines the parts of \method{} in turn: the renewal of the
pairs and the reference model, the two terms of the loss, the warm-up
rounds and the server average. Table~\ref{tab:ablation} reports the
variants, each tested against \method{} at the same number of rounds:
three, that is, a single preference round, in the lower block, and four
in the upper. All but the two with another number of warm-up rounds
start from the model of the first row, FedAvg SFT after
\setting{rounds.warmup} rounds. Significance is stated after the
correction of Section~\ref{sec:stat-test}; a $p$-value is unadjusted
unless marked as adjusted.

\begin{table}[t]
\caption{EX (\%) of \method{} and of its variants on BIRD dev and Spider
test; the first row is the model after the \setting{rounds.warmup}
warm-up rounds. $\Delta$ is the difference in EX to \method{} of the
same block, and $p$ the unadjusted $p$-value of the exact McNemar test
against that row, with an asterisk where the difference is significant
at the \fact{test.alpha} level after the Holm correction over the
comparisons of the table. ``Clients pooled'' is run on Spider
only}\label{tab:ablation}
\centering
%% Written by tools/paper_numbers.py. Do not edit.
\setlength{\tabcolsep}{3pt}%
\small
\begin{tabular}{@{}lrrlrrlrrl@{}}
\toprule
 & \multicolumn{6}{c}{BIRD dev} & \multicolumn{3}{c}{Spider test} \\
\cmidrule(lr){2-7}\cmidrule(lr){8-10}
 & EX & $\Delta$ & $p$ & TS & $\Delta$ & $p$ & EX & $\Delta$ & $p$ \\
\midrule
FedAvg SFT, 2 rounds & 42.70 &  &  & 36.11 &  &  & 74.10 &  &  \\
\midrule
\multicolumn{10}{@{}l}{\textit{4 rounds}} \\
\midrule
\method & 48.76 &  &  & 41.59 &  &  & 78.20 &  &  \\
FedAvg SFT & 42.57 & $-$6.19 & $1.8\times10^{-10}$\textsuperscript{*} & 36.70 & $-$4.89 & $7.9\times10^{-7}$\textsuperscript{*} & 74.71 & $-$3.49 & $1.9\times10^{-6}$\textsuperscript{*} \\
Fixed pairs and reference & 46.15 & $-$2.61 & 0.0091\textsuperscript{*} & 39.05 & $-$2.54 & 0.0093\textsuperscript{*} & 78.02 & $-$0.18 & 1.00 \\
No NLL term & -- & -- & -- & -- & -- & -- &  &  &  \\
No preference term & 45.24 & $-$3.52 & $9.3\times10^{-4}$\textsuperscript{*} & 38.20 & $-$3.39 & 0.0015\textsuperscript{*} & -- & -- & -- \\
No warm-up round & -- & -- & -- & -- & -- & -- & 78.02 & $-$0.18 & 1.00 \\
Three warm-up rounds & -- & -- & -- & -- & -- & -- & 77.97 & $-$0.23 & 1.00 \\
Clients pooled &  &  &  &  &  &  & -- & -- & -- \\
\midrule
\multicolumn{10}{@{}l}{\textit{3 rounds}} \\
\midrule
\method & 46.22 &  &  & 39.05 &  &  & 77.46 &  &  \\
No NLL term & 44.98 & $-$1.24 & 1.00 & 38.92 & $-$0.13 & 1.00 & 73.59 & $-$3.87 & $1.5\times10^{-7}$\textsuperscript{*} \\
No preference term & 44.39 & $-$1.83 & 0.74 & 37.35 & $-$1.70 & 0.75 & 75.87 & $-$1.59 & 0.21 \\
Averaged by questions & 46.22 & +0.00 & 1.00 & 39.31 & +0.26 & 1.00 & 77.78 & +0.32 & 1.00 \\
Averaged equally & 46.02 & $-$0.20 & 1.00 & 38.92 & $-$0.13 & 1.00 & 76.94 & $-$0.52 & 0.98 \\
\bottomrule
\end{tabular}

\end{table}

With the pairs and the reference model held fixed, EX is lower on BIRD
dev and level on Spider test. In the row ``Fixed pairs and reference'',
the second preference round trains on the pairs of the first preference
round, against the reference model of that round; the NLL term is
retained. FedDPO (Section~\ref{sec:baselines}) and Liu et
al.~\cite{liu2025cotdpo} likewise keep one set of pairs and one
reference model throughout training. The row fixes both together and
does not show from which of the two a difference arises. On BIRD dev, it
reaches \ex{bird}{fixed}{4}, against \ex{bird}{k2}{4} for \method{}
after four rounds, a significant difference
($p =$~\pval{bird}{k2}{4}{fixed}{4}), and \ex{bird}{k2}{3} for \method{}
after three. On Spider test, the two do not differ
(\ex{spider}{fixed}{4} against \ex{spider}{k2}{4},
$p =$~\pval{spider}{k2}{4}{fixed}{4}); there, the second preference
round of \method{} itself adds \gap{spider}{k2}{3}{k2}{4} points,
against \gap{bird}{k2}{3}{k2}{4} on BIRD dev (Section~\ref{sec:rq1}).

Without the NLL term, a preference round brings no gain on Spider test.
The row ``Without the NLL term'' is one preference round with the
preference term alone, the loss of FedDPO and of Liu et
al.~\cite{liu2025cotdpo}; in a single round, this is its only difference
from \method. On Spider test, the row reaches \ex{spider}{feddpo}{3},
against \ex{spider}{fedavg}{2} for the model before the round, and is
significantly lower than \method, by \gap{spider}{k2}{3}{feddpo}{3}
points ($p =$~\pval{spider}{k2}{3}{feddpo}{3}). On BIRD dev, it reaches
\ex{bird}{feddpo}{3}, against \ex{bird}{fedavg}{2} before the round,
\gain{bird}{k2}{3}{feddpo}{3} points relative to \method{}
($p =$~\pval{bird}{k2}{3}{feddpo}{3}).
CodeLutra~\cite{tao2024codelutra} on Spider and Pang et
al.~\cite{pang2024irpo} on GSM8K likewise report a lower accuracy with
the DPO term alone than with both terms.

With the preference term, fewer predictions end in an execution error; a
difference in EX is not established on either dataset. The row ``Without
the preference term'' trains with the NLL term alone on the correct
sample of each pair, in an otherwise unchanged preference round. After
one such round, its EX is \gap{bird}{k2}{3}{rft}{3} points below that of
\method{} on BIRD dev and \gap{spider}{k2}{3}{rft}{3} on Spider test,
neither difference being significant (adjusted
$p =$~\padj{bird}{k2}{3}{rft}{3} and \padj{spider}{k2}{3}{rft}{3}). It
leaves \nerr{bird}{rft}{3} execution errors on BIRD dev and
\nerr{spider}{rft}{3} on Spider test, against \nerr{bird}{k2}{3} and
\nerr{spider}{k2}{3} for \method{} and \nerr{bird}{fedavg}{2} and
\nerr{spider}{fedavg}{2} before the round, both differences to \method{}
being significant under the test of Section~\ref{sec:stat-test} applied
to execution errors ($p =$~\errpval{bird}{k2}{3}{rft}{3} and
\errpval{spider}{k2}{3}{rft}{3}). \method{} has more incorrect results,
\nwr{bird}{k2}{3} against \nwr{bird}{rft}{3} on BIRD dev
($p =$~\wrpval{bird}{k2}{3}{rft}{3} under the same test, significant)
and \nwr{spider}{k2}{3} against \nwr{spider}{rft}{3} on Spider test
($p =$~\wrpval{spider}{k2}{3}{rft}{3}, not significant).

At the same total of four rounds, EX on Spider test is level with no,
two and three warm-up rounds (\ex{spider}{k0}{4}, \ex{spider}{k2}{4} and
\ex{spider}{k3}{4}; neither difference to \method{} is significant);
what differs is test-suite accuracy and the number of training questions
without a correct sample. Without warm-up rounds, the model is above
FedAvg SFT in EX on Spider dev, Spider-Syn and Spider-Realistic and
below it in TS on all three (\ts{spiderdev}{k0}{4} against
\ts{spiderdev}{fedavg}{4} on Spider dev; Table~\ref{tab:robust}),
whereas \method{} is above FedAvg SFT in both metrics
(Section~\ref{sec:rq3}). This is a comparison of values; the tested one
is with \method, against which the lower TS of this variant is
significant on Spider-Syn only (\tsgap{syn}{k0}{4}{k2}{4} points,
$p =$~\tspval{syn}{k0}{4}{k2}{4}). On Spider-DK, the EX of the variant
is \gap{dk}{k0}{4}{k2}{4} points above that of \method, which is not significant (adjusted
$p =$~\padj{dk}{k0}{4}{k2}{4}).

In training, the warm-up rounds leave fewer questions without a correct
sample, that is, with none of the $N$ samples of the first preference
round labeled correct (Section~\ref{sec:warmup}). Of the
\fact{cost.bird.pref1.questions} training questions of BIRD,
\fact{warmup.bird.k0.no_correct} are such questions with samples from
the pretrained model and \fact{warmup.bird.k2.no_correct} after two
warm-up rounds; of the \fact{train.spider.with_rows} questions of Spider
on databases with rows, \fact{warmup.spider.k0.no_correct},
\fact{warmup.spider.k2.no_correct} and
\fact{warmup.spider.k3.no_correct} with no, two and three warm-up
rounds. \method{} learns from the gold SQL query of such a question in
the warm-up rounds; an alternative also pairs it, as the correct side,
with $M$ incorrect samples of the question, drawn as in
Section~\ref{sec:pref-round}. We evaluate it on BIRD in the variant
without warm-up rounds, where these questions are most numerous, in one
preference round on the same samples and labels. With these pairs, EX on
BIRD dev is \gap{bird}{k0}{1}{k0gold}{1} points lower
(\ex{bird}{k0gold}{1} against \ex{bird}{k0}{1}), which is not
significant ($p =$~\pval{bird}{k0}{1}{k0gold}{1}; adjusted,
\padj{bird}{k0}{1}{k0gold}{1}). Execution errors fall from
\nerr{bird}{k0}{1} to \nerr{bird}{k0gold}{1}, and incorrect results rise
from \nwr{bird}{k0}{1} to \nwr{bird}{k0gold}{1}.

Averaging the adapters of the first preference round again with other
weights leads to no significant difference in EX. With the question
counts $n_k$ as weights, the weight of a client differs from its weight
by pair counts by at most \fact{weights.bird.questions.max_shift}
percentage points on BIRD and \fact{weights.spider.questions.max_shift}
on Spider; with equal weights, by at most
\fact{weights.bird.equal.max_shift} and
\fact{weights.spider.equal.max_shift}. The largest difference in EX to
\method{} is \gap{spider}{k2}{3}{reagge}{3} points, for equal weights on Spider test
($p =$~\pval{spider}{k2}{3}{reagge}{3}).

The row ``Clients pooled'' changes the setting and not the method: it
shows whether the preference rounds lose by running at five clients with
a server average. Its preference rounds run on the training data of the
five clients held at one site, as a single client. On Spider test, it
reaches \ex{spider}{pooled}{4},
\gain{spider}{k2}{4}{pooled}{4} points relative to \method{}
($p =$~\pval{spider}{k2}{4}{pooled}{4}).
